\section{Hereditary caps and graph notation}
\label{sec:setup}

For integers $k\geq1$ and $n\geq0$, write $n=kq+s$ with $0\leq s<k$ and define
\[
\tau_k(n)=k\binom{q}{2}+qs.
\]
This is the minimum number of edges in a graph on $n$ vertices with independence number at most $k$, by applying Tur\'an's theorem \cite{Turan1941} to the complement. Equivalently, $\binom{n}{2}-\tau_k(n)$ is the Tur\'an number $\ex(n,K_{k+1})$.

Define the target cap
\[
D(n)=\binom{n}{2}-11\tau_{12}(n).
\]
The values used in the proof are
\[
\begin{array}{c|rrrrrrrr}
n&13&14&23&24&36&37&38&49\\ \hline
D(n)&67&69&132&144&234&237&241&\text{not used}
\end{array}
\]
(the $49$-vertex contradiction below compares a target lower bound with a different color-class upper bound).

\begin{lemma}[Hereditary color caps]
\label{lem:caps}
In a balanced $12$-edge-coloring of $K_{145}$, every color graph $C_i$ satisfies $\alpha(C_i)\leq12$. Consequently, for every vertex set $U$,
\[
e(C_i[U])\geq\tau_{12}(|U|).
\]
If $G$ is a color graph with the fewest edges, then
\[
e(G)\leq870,\qquad
e(G[U])\leq D(|U|)\quad\text{for every }U\subseteq V(K_{145}),
\]
and $\alpha(G)\leq12$, $\omega(G)\leq12$.
\end{lemma}

\begin{proof}
An independent set of size $13$ in a color graph would be a $13$-set containing no edge of that color, contrary to balance. Thus $\alpha(C_i)\leq12$. For any $U$, the complement of $C_i[U]$ is $K_{13}$-free, so Tur\'an's theorem gives
\[
e(\overline{C_i[U]})\leq \ex(|U|,K_{13})
=\binom{|U|}{2}-\tau_{12}(|U|).
\]
Subtracting from $\binom{|U|}{2}$ gives the lower bound on $e(C_i[U])$. Since $G$ is a least-edge color class, $12e(G)\leq\binom{145}{2}=10440$, so $e(G)\leq870$. On any $U$, the eleven non-target colors contribute at least $11\tau_{12}(|U|)$ edges, leaving at most $D(|U|)$ target edges. Finally, $\alpha(G)\leq12$ was already shown, and a $K_{13}$ in $G$ would be a $13$-set containing no edge of any other color, also impossible.
\end{proof}

\begin{definition}
For $a,n\geq0$, let $\mathcal R(a,n)$ be the class of simple graphs $J$ on $n$ vertices such that
\[
\alpha(J)\leq a,\qquad \omega(J)\leq11,\qquad
e(J[U])\leq D(|U|)\quad\text{for every }U\subseteq V(J).
\]
The class is hereditary under induced subgraphs.
\end{definition}


The hereditary condition in $\mathcal R(a,n)$ is essential in the lower-order arguments. In contrast, the abstract $47$-vertex low-edge and equality theorems below use only their five explicitly displayed graph hypotheses and do not assume the caps $D$.

We also use the following second-Tur\'an theorem of Kang and Pikhurko \cite{KangPikhurko2005}. If $n\geq r+3$ and $2\leq r\leq(n-1)/2$, then every $K_{r+1}$-free graph $L$ on $n$ vertices that is not $r$-partite satisfies
\[
e(L)\leq t_r(n)-\left\lfloor\frac nr\right\rfloor+1,
\qquad t_r(n)=\ex(n,K_{r+1}).
\]
Our applications are $(n,r)=(25,2)$ and $(49,12)$, both within the stated parameter range.
